\begin{center}
  \large\textbf{ABSTRACT}\\[2ex]
  \large\textbf{\engtatitle{}}
\end{center}

\addcontentsline{toc}{chapter}{ABSTRAK}

\vspace{2ex}

\begingroup
% Menghilangkan padding
\setlength{\tabcolsep}{0pt}

\noindent
\begin{tabularx}{\textwidth}{l >{\centering}m{2em} X}
  Name / NRP   & : & \name{} / \nrp{}        \\

  Department & : & \engdepartment{} - Faculty of \engfaculty     \\

  Advisor        & : & 1. \advisor{}   \\
                    &   & 2. \coadvisor{} \\
\end{tabularx}
\endgroup

Manual student attendance tracking in university environments is often inefficient and prone to fraud. Although automated facial recognition-based attendance systems are widely implemented, conventional \textit{Centralized Learning} methods raise serious concerns by requiring raw student face images to be sent to a centralized server. This triggers high risks of sensitive biometric privacy leaks and excessive network bandwidth consumption. To address these challenges, this study proposes a distributed attendance system based on \textit{Edge Computing} that integrates the \textit{Federated Learning} method with the \textit{MobileFaceNet} model architecture and the \textit{FedProx} optimization algorithm. All system components are packaged using \textit{Docker} containerization to be deployed on a heterogeneous edge hardware cluster consisting of Jetson Nano and Raspberry Pi 3 Model B. This approach facilitates local and collaborative model training on edge devices without transmitting raw biometric data out of the respective devices. System evaluation covers training metrics, network bandwidth efficiency, and inference reliability under real-life and video simulation scenarios. Experimental results show that the \textit{Federated Learning} global model achieves an accuracy of 98.16\% with a validation loss of 0.3599, comparable to the centralized method which records a 96.93\% accuracy. Furthermore, this distributed method significantly saves bandwidth usage by only transmitting 159.46 MB of parameter data compared to the centralized method's 1152.34 MB. In real-life inference testing, the system achieves a 100\% operational success rate at distances of 0.5 to 3 meters under both sufficient and low lighting conditions, maintaining a stable classification accuracy above 90\% using a similarity threshold of 0.7. Local processing latency on Jetson Nano is highly responsive ($\sim$700 ms), whereas Raspberry Pi requires $\sim$3100 ms due to RAM limitations. This study concludes that the integration of \textit{Federated Learning} and \textit{Edge Computing} is proven effective in delivering a secure, responsive, and bandwidth-efficient attendance system.

% Ubah kata-kata berikut dengan kata kunci dari tugas akhir
\textbf{Keywords: \keywords{}}
